\documentclass[12pt]{article}
\usepackage[T1]{fontenc}
\usepackage[utf8]{inputenc}
\usepackage[brazil]{babel}
\usepackage[margin=1in]{geometry}
\setlength{\parindent}{0pt}
\begin{document}
\section*{Tema 94}
Processo: PEDILEF 2005.37.00.749443-3/ MA\\
Situacao: Julgado\\
Ramo: DIREITO PREVIDENCIÁRIO\\
Questao: Saber se o rito da Lei n. 10.259/2001 pode ser aplicado no âmbito dos juizados especiais estaduais para julgamento das ações previdenciárias, em razão da competência delegada (CF/88 art. 109, § 3º).\\
Tese: É absoluta a incompetência do Juizado Especial Cível Estadual para o processamento e julgamento das causas previdenciárias, por expressa vedação legal à aplicação da Lei n. 10.259/2001 no âmbito do juízo estadual.\\
Fonte: https://www.cjf.jus.br/phpdoc/virtus/
\end{document}
